\documentclass[11pt]{article}
\usepackage[a4paper,margin=25mm]{geometry}
\usepackage{amsmath,amssymb,amsthm}
\usepackage[hidelinks]{hyperref}
\emergencystretch=2em
\newtheorem{theorem}{Theorem}
\title{Two Positive Roots in a Geometric Lattice:\\A Disproof of Conjecture 00000008551}
\author{GitHub: gaochengzhecpu}
\date{}
\begin{document}
\maketitle
\begin{abstract}
The partition lattice of a three-element set is geometric and has
characteristic polynomial \mbox{$(t-1)(t-2)$}. Its two distinct roots are positive, so their
signs do not alternate. The associated triangle graph has chromatic
polynomial $t(t-1)(t-2)$, with the same two consecutive nonzero roots.
We give the actual lattice, its rank and M\"obius recurrence, and a
counting proof for the graph polynomial. Lean verifies these objects
and refutes the corresponding universally quantified root-sign clauses.
The scope of the source's nonstandard terminology is stated explicitly.
\end{abstract}

\section*{Reading of the bilingual source}
Both languages assert that the \emph{roots} of the chromatic polynomial
of a geometric lattice alternate in sign. They do not say that its
\emph{coefficients} alternate in sign. The source's preliminary phrase
``intersection-disjoint maps'' does not define a class of maps, a color
set, or a counting rule. Accordingly we state precisely the two standard
interpretations covered here.

For a finite geometric lattice $L$ of rank $r$ with rank function $\rho$,
the standard lattice characteristic polynomial is
\[
 \chi_L(t)=\sum_{x\in L}\mu(0,x)t^{r-\rho(x)},
 \qquad
 \sum_{y\leq x}\mu(0,y)=
 \begin{cases}1&x=0,\\0&x\ne0.\end{cases}
\]
This is the usual lattice extension of a graph chromatic polynomial.
For a graph $G$, its bond lattice $L_G$ consists of the partitions of
its vertices into connected blocks, ordered by refinement. The usual
graph polynomial $P_G(t)$ counts proper colorings and satisfies
$P_G(t)=t^{c(G)}\chi_{L_G}(t)$. These conventions and their relation
are recorded in Stanley's publicly available \emph{Enumerative
Combinatorics}, Volume 1, second-edition manuscript, equation (3.39)
and Exercise 3.108(b)--(c)~\cite{stanley}.

We disprove the root assertion with either of these standard
normalizations. We do not assign a new meaning to the undefined phrase
in the source. If a different polynomial is intended, its definition
would first have to be supplied. The true coefficient-sign theorem is
not refuted: our polynomial has coefficients $1,-3,2$, which do
alternate. Neither version of the separate root-modulus assertion is
needed to refute the conjunction.

\section*{The geometric lattice and its polynomial}
Let $L=\Pi_3$ be the lattice of all partitions of $\{0,1,2\}$, ordered
by refinement. Its five elements are
\[
 \begin{array}{c|c|c}
 \text{element}&\text{partition}&\rho\\\hline
 0&0\mid1\mid2&0\\
 a&01\mid2&1\\
 b&02\mid1&1\\
 c&0\mid12&1\\
 1&012&2
 \end{array}
\]
The elements $a,b,c$ are pairwise incomparable. Two distinct such
elements have meet $0$ and join $1$. These facts, together with the
least and greatest elements, specify all meets, joins and order
relations.

Every element is a join of atoms: the empty join gives $0$, each of
$a,b,c$ is an atom, and any two distinct atoms join to $1$. The lattice
is modular. Indeed, the modular law for $x\leq z$ is immediate if
$x=0$, $z=1$, or $x=z$; these exhaust its comparable pairs. Modularity
implies upper semimodularity. Thus this finite atomistic semimodular
lattice is geometric. The displayed rank starts at zero and rises by
one at each covering relation.

The bottom-row M\"obius recurrence determines
\[
 \mu(0,0)=1,\qquad
 \mu(0,a)=\mu(0,b)=\mu(0,c)=-1,
 \qquad \mu(0,1)=2.
\]
For an atom the recurrence is $1+\mu(0,a)=0$; at the top it is
$1-1-1-1+\mu(0,1)=0$. These equations also show uniqueness of
the listed values, rather than merely stipulating five numbers.
Consequently
\[
 \chi_L(t)=t^2-3t+2=(t-1)(t-2).
\]

\begin{theorem}
The roots of the characteristic polynomial of a finite geometric
lattice need not alternate in sign. The same failure occurs for the
graph chromatic polynomial associated with its bond lattice.
\end{theorem}
\begin{proof}
For $L=\Pi_3$, the two distinct roots of $\chi_L$ are $1$ and $2$.
They are simple and consecutive and their product is $2>0$.
Thus they have the same sign. This obstruction is unaffected by
whether roots are ordered increasingly, decreasingly, or by modulus.

For the second assertion take the actual complete graph $K_3$ on
$\{0,1,2\}$. Every nonempty subset induces a connected graph, so
its bond lattice is exactly $\Pi_3$. A proper coloring with $q$ labeled
colors is an injection from its three vertices into the color set.
There are $q(q-1)(q-2)$ such injections for $q\geq2$, and zero
for $q=0,1$; the same polynomial gives zero in those cases. Therefore
\[
 P_{K_3}(t)=t(t-1)(t-2)=t\chi_L(t).
\]
This identification is a counting identity for every nonnegative
integer $q$, which uniquely determines the polynomial. Its roots
$1$ and $2$ are again consecutive, distinct, nonzero, and positive.
Ignoring the additional root $0$ cannot repair sign alternation.
\end{proof}

\section*{Formal verification and scope}
The Lean type \texttt{ThreePartitions} carries a proved Mathlib
\texttt{Lattice}, bounded order, modularity, and
\texttt{IsAtomistic} instance. Its rank function is checked using
Mathlib's covering relation. The explicit block-label functions define
actual equivalence relations on \texttt{Fin 3}. The order isomorphism
\begin{center}\small\texttt{partitionOrderIso}\end{center}
identifies the five-element construction with \emph{all} elements of
\texttt{Setoid (Fin 3)}, ordered by inclusion of equivalence relations.
Thus the lattice is connected to the standard partition object, and
not just to a numerical rank list.

Lean proves that the supplied integer function satisfies the finite
bottom-row recurrence and is its unique solution on this lattice.
The recurrence is named \texttt{IsMobiusRow} in the source.
No general incidence-algebra package is imported.
The polynomial is the actual finite sum over the lattice of M\"obius
values times rank monomials; its factorization and its complete real
root set are proved.

\begin{samepage}
The proposition \texttt{RootSignClause} quantifies over all finite
bounded atomistic upper-semimodular lattices, all rank functions, and
all bottom-row M\"obius functions. It asserts the following necessary
condition of alternating root signs: any two consecutive distinct
nonzero real roots have negative product. Its negation is
\begin{center}\small\texttt{disprove\_root\_sign\_clause}.\end{center}
\end{samepage}
This deliberately weak necessary condition already fails at $1,2$;
no issue involving repeated, zero, or nonreal roots is used.

For the graph interpretation, Lean uses Mathlib's actual complete
\texttt{SimpleGraph (Fin 3)} and \texttt{Coloring}. It constructs an
equivalence with embeddings of \texttt{Fin 3} into \texttt{Fin q}
and proves the coloring count for \emph{every} $q$. The polynomial
\texttt{triangleChromatic} is proved to evaluate to that count.
\texttt{GraphRootSignClause} quantifies over finite graphs and
polynomials having this counting property; its negation is
\begin{center}\small\texttt{disprove\_graph\_root\_sign\_clause}.\end{center}
The elementary identification of the bond lattice of $K_3$ with
$\Pi_3$, and the general graph--bond-lattice formula quoted for
context, are proved or explained on paper rather than formalized as
general graph-connectivity theorems. Both specific polynomials and
both counterexample statements are formalized.

The Lean proofs use no \texttt{sorry}, \texttt{admit},
\texttt{native\_decide}, or custom axiom. Their axiom reports contain
only \texttt{propext}, \texttt{Classical.choice}, and
\texttt{Quot.sound}. Finite kernel checks certify the concrete lattice
axioms and recurrence; the universal claims are explicitly quantified
and refuted by the resulting objects.

\section*{Source and reproduction}
The exact bilingual source is included as \texttt{SOURCE.md}, copied
from the upstream file
\begin{center}\small\texttt{conjectures/00000008551.md}.\end{center}
The provenance record gives the upstream revision and source hash.
The project pins Lean 4.19.0 and Mathlib revision
\begin{center}\small\texttt{c44e0c8ee63ca166450922a373c7409c5d26b00b}.\end{center}
From \texttt{lean/}, with the pinned dependencies available, run
\texttt{lake build} and
\begin{center}\small\texttt{lake env lean -DwarningAsError=true Main.lean}.\end{center}
Compile this document with Tectonic or a compatible LaTeX engine.
The verification directory records author checks and subsequent
coordinator checks separately. This draft was developed by a delegated
AI agent with a self-review; the coordinating agent reviews it
separately. No external independent review is claimed.

\begin{thebibliography}{1}
\bibitem{stanley} R. P. Stanley, \emph{Enumerative Combinatorics},
Volume 1, second-edition manuscript, version of 15 July 2011.
Section 3.3 (geometric lattices), equations (3.38)--(3.39),
and Exercise 3.108(b)--(c).
\href{https://math.mit.edu/~rstan/ec/ec1.pdf}{Author's public manuscript}.
\end{thebibliography}
\end{document}
